% Related work, drafted with an assistant and not yet edited.
% The fixture for the paper branch: the markup section does NOT apply here.

\section{Related Work}
\label{sec:related}

Reinforcement learning has been applied to adaptive traffic signal control since
at least 2016~\cite{genders2016,wei2019}, and the actor-critic designs that
followed build directly on the earlier value-based methods~\cite{mnih2015}.

Tabular $Q$-learning remains the canonical starting point. It is interpretable
in a way deep methods are not, which is what makes it useful for baselines on
small intersection topologies~\cite{abdulhai2003}. Function approximation came
next and generalised across unseen states, which cut the number of episodes
needed to reach a fixed return. The resulting methods converged stably on
networks of 4--9 intersections, provided the reward was shaped carefully.

Queue length and phase stability interact, and that interaction has received
little attention: prior studies optimise average delay in isolation. This work
measures the two together---the distinction the rest of the paper turns on.

The literature covers the single-intersection case well and does not scale to
corridors. Coordination overhead may dominate at higher flow densities, though
only two studies have measured it, and both report the effect to be robust to
sensor dropout~\cite{chu2020}. Coordination across corridors therefore remains
an open problem, and Section~\ref{sec:method} sets out the two-branch design we
use to attack it.
